\chapter{Kiểm thử, triển khai và vận hành}

\section{Chiến lược kiểm thử}

\begin{longtable}{L{3.4cm} L{2.6cm} L{7.6cm}}
\toprule
\textbf{Phạm vi} & \textbf{Số tệp} & \textbf{Lệnh và nội dung}\\
\midrule
\endhead
Backend & 102 & \code{pytest} từ \code{backend-service/} — kiểm thử đơn vị dịch vụ, kiểm thử tuyến API với cơ sở dữ liệu SQLite bất đồng bộ.\\
AI service & 61 & \codesp{pytest tests/} — trong đó \code{tests/trace\_cag/} là kiểm thử tích hợp pipeline, \code{tests/content\_etl/} kiểm thử bộ tiếp hợp ETL.\\
Flutter & 120 & \codesp{flutter test} từ \code{flutter-app/} — widget, provider, use case.\\
Liên dịch vụ & — & \code{system-testing/} và \code{scripts/smoke-prod.sh}.\\
\bottomrule
\caption{Bộ kiểm thử}
\end{longtable}

\subsection{Quy tắc về giả lập}

\begin{luuy}[Bản giả không được dễ dãi hơn bản thật]
Phải dùng \code{autospec=True} khi vá hàm. Một \code{AsyncMock} trần đã chấp nhận tham số
\code{httpx\_module=} không còn tồn tại suốt nhiều tháng, trong khi
\code{/api/v1/ai/translate} trả về rỗng cho mọi từ. Kiểm thử vẫn xanh vì bản giả nhận mọi
tham số; chỉ có người dùng thật mới phát hiện ra.
\end{luuy}

Tương tự với dữ liệu mẫu: tệp mẫu của ETL phải có \textbf{cùng hình dạng} với tệp tải về
thật (nén, mã hoá ký tự, \en{namespace}), nếu không bộ tiếp hợp sẽ vượt kiểm thử nhưng
hỏng ở sản xuất.

\subsection{Kiểm thử bảo vệ bất biến}

Một số kiểm thử tồn tại không phải để kiểm tra tính năng mà để chặn hồi quy kiến trúc:

\begin{itemize}[leftmargin=1.4em]
  \item \code{test\_content\_contract\_parity.py} — bảo đảm danh sách khoá được phép của
  quan sát giống nhau ở cả hai dịch vụ.
  \item \code{test\_container\_hardening.py} — bảo đảm container không chạy bằng
  \code{root} và không mở cổng thừa.
  \item \code{test\_production\_kg\_isolation.py} — bảo đảm đồ thị tri thức sản xuất không
  bị kiểm thử ghi đè.
  \item \code{test\_dependency\_trace.py} — bảo đảm vết phụ thuộc của bộ đệm được biên dịch
  đúng, vì sai vết nghĩa là phản hồi cũ không bị vô hiệu.
\end{itemize}

\section{Đóng gói và triển khai}

\subsection{Cấu hình sản xuất}

Sản xuất chạy trên một máy chủ ảo tại \code{/opt/lexilingo} (bí danh
\code{sgp1-01-lexi}), nhánh \code{dev}, khởi động bằng Docker Compose với hai tệp môi
trường: \code{.env.production} và \code{.env.production.secrets}.

Các dịch vụ trong \code{docker-compose.yml}: \code{gateway}, \code{postgres},
\code{mongodb}, \code{redis}, \code{backend-service},
\code{backend-reminder-worker}, \code{backend-reminder-beat}, \code{ai-service},
\code{prometheus}, \code{grafana}.

\subsection{Quy trình triển khai chuẩn}

\begin{lstlisting}[language=bash]
./scripts/deploy-one-shot.sh
# pull -> up -d --remove-orphans -> cho tung dich vu khoe -> smoke test
\end{lstlisting}

Compose chỉ dựng lại container nào thực sự thay đổi và giữ container cũ chạy cho tới khi
container mới sẵn sàng.

\subsection{Ba điều cấm tuyệt đối}

Ba sai lầm dưới đây đã làm sập ai-service ngày 14/8/2026 và phá huỷ ảnh dùng để quay lui.
Mỗi điều bị cấm độc lập:

\begin{longtable}{L{3.6cm} L{5.0cm} L{5.0cm}}
\toprule
\textbf{Không bao giờ} & \textbf{Vì sao} & \textbf{Thay bằng}\\
\midrule
\endhead
\code{--force-recreate} & Phá container đang chạy \textbf{trước khi} container thay thế được chứng minh là hoạt động. Khi bản mới không khởi động được, không còn gì phục vụ. & \code{up -d} thuần; compose tự dựng lại khi có thay đổi thật.\\
\codesp{timeout N} bọc quanh lệnh dựng ảnh & Bánh xe \code{llama-cpp-python} mất 15--30 phút trên 2 vCPU. \code{timeout 240} huỷ giữa chừng, để lại thẻ \code{:latest} trỏ vào hư không — ảnh quay lui biến mất. & Không đặt thời hạn; chạy tách rời rồi theo dõi nhật ký.\\
Dựng ảnh trong lệnh SSH tiền cảnh & Ống dữ liệu đứt giữa chừng giết tiến trình từ xa; ảnh dựng xong nhưng container không được đổi. & \codesp{setsid nohup ... > /tmp/deploy.log 2>\&1 < /dev/null \&} rồi theo dõi tệp nhật ký.\\
\bottomrule
\caption{Ba thao tác bị cấm khi triển khai}
\end{longtable}

Ngoài ra: \code{| tail -N} khi dựng ảnh sẽ giấu toàn bộ tiến trình cho tới khi kết thúc —
hãy ghi ra tệp rồi theo dõi tệp.

\subsection{Luôn giữ một lối thoát}

Không có bản sao ảnh và không có kho ảnh riêng, nên một thẻ bị mất là mất vĩnh viễn.
Trước khi động vào dịch vụ đang chạy:

\begin{lstlisting}[language=bash]
sudo docker tag lexilingo-ai-service:latest lexilingo-ai-service:rollback
\end{lstlisting}

Và luôn tách dựng ảnh khỏi đổi container, để một lần đổi hỏng vẫn còn container cũ:

\begin{lstlisting}[language=bash]
$COMPOSE build ai-service   # container cu van dang phuc vu
$COMPOSE up -d ai-service   # chi doi sau khi anh da ton tai
\end{lstlisting}

\subsection{Mức độ ảnh hưởng khi mất từng dịch vụ}

Mất ai-service làm mất hội thoại, dịch, nhận dạng giọng nói và TRACE-CAG, \textbf{nhưng
không} làm mất khoá học, bài học, đăng nhập hay tiến độ — những thứ đó thuộc
backend-service. Thời gian dựng lại ai-service là 15--30 phút, backend-service khoảng
1 phút. Vì vậy khi phải chọn thứ tự triển khai, backend đi trước.

\section{Quan trắc}

\begin{itemize}[leftmargin=1.4em]
  \item \textbf{Chỉ số}: Prometheus thu thập từ cả hai dịch vụ; Grafana hiển thị. Bảng
  điều khiển gồm độ trễ theo phân vị, tỉ lệ lỗi, tỉ lệ trúng bộ đệm, mức dùng bộ nhớ mô
  hình.
  \item \textbf{Nhật ký}: có cấu trúc, mang \code{X-Request-ID}; tác vụ Celery mang
  \code{task:\{id\}} nên chuỗi log của một tác vụ nền cũng gom được thành một khối.
  \item \textbf{Vết}: \code{telemetry.py} phía AI service dùng OpenTelemetry.
  \item \textbf{Kiểm tra sức khoẻ}: \code{/health} (còn sống) và \code{/health/ready}
  (sẵn sàng nhận lưu lượng) — phân biệt này quan trọng vì ai-service cần thời gian nạp mô
  hình trước khi thực sự phục vụ được.
  \item \textbf{Cảnh báo bảo mật}: \code{scripts/gateway-security-alerts.sh} và
  \code{scripts/lexilingo-sentinel.py}.
\end{itemize}

\section{Quy trình phát triển}

\begin{longtable}{L{4.0cm} L{9.6cm}}
\toprule
\textbf{Bước} & \textbf{Lệnh}\\
\midrule
\endhead
Môi trường backend & \codesp{cd backend-service \&\& source venv/bin/activate}\\
Môi trường AI & \codesp{cd ai-service \&\& source venv/bin/activate} (Python 3.12)\\
Di trú cơ sở dữ liệu & \codesp{alembic upgrade head}\\
Chạy backend & \codesp{uvicorn app.main:app --reload --port 8000}\\
Chạy toàn cục bộ & \code{./scripts/start-all.sh}, dừng bằng \code{./scripts/stop-all.sh}\\
Kiểm tra ngăn xếp cục bộ & \code{./scripts/validate\_local\_stack.sh}\\
\bottomrule
\caption{Lệnh phát triển thường dùng}
\end{longtable}

Quy trình CI trên GitHub Actions: chạy kiểm thử cho mỗi \en{pull request}, triển khai khi
gộp vào nhánh chính, đồng bộ chuỗi dịch với Crowdin, rà soát mã bằng AI, và tự động gộp
cập nhật phụ thuộc.

\section{Hạn chế hiện tại và hướng phát triển}

\subsection{Hạn chế đã biết}

\begin{enumerate}[leftmargin=1.4em]
  \item \textbf{Sao lưu chưa ra khỏi máy chủ.} Bản sao lưu nằm cùng ổ đĩa với cơ sở dữ
  liệu; cần cấu hình \code{OFFSITE\_RCLONE\_REMOTE}.
  \item \textbf{Giới hạn của pipeline quan sát chưa được đo tải thật.} Các con số hiện
  hành là ranh giới an toàn, chưa phải năng lực đã kiểm chứng.
  \item \textbf{Hình dạng bài tập không được kiểm tra khi ghi.} Một bài tập méo mó chỉ bị
  phát hiện bởi người học không trả lời được.
  \item \textbf{Mười ba khoá học lớn không tái tạo được từ kho mã.}
  \item \textbf{Tiêu thụ thụ động chưa tạo bằng chứng học tập} — đúng theo thiết kế, nhưng
  đồng nghĩa với việc phần lớn thời gian dùng ứng dụng ở các mục nội dung không được đo.
  \item \textbf{Đường sinh nội dung im lặng rơi về khuôn mẫu} khi LLM lỗi, khiến sự cố khó
  phát hiện.
\end{enumerate}

\subsection{Hướng phát triển}

\begin{enumerate}[leftmargin=1.4em]
  \item Bổ sung bài chính tả hoặc quiz cho podcast và YouTube để nối chúng vào hệ thống
  ghi nhận.
  \item Kiểm tra hình dạng bài tập ngay tại thời điểm ghi, thay vì chỉ kiểm toán sau.
  \item Phát tín hiệu rõ ràng khi tác tử nội dung rơi về khuôn mẫu, thay vì im lặng.
  \item Hoàn tất đợt thử tải theo giai đoạn cho pipeline quan sát và ghi kết quả vào
  \code{docs/operations/learner-state-rollout.md}.
  \item Đưa sao lưu ra ngoài máy chủ.
  \item Mở rộng đánh giá TRACE-CAG: so sánh có kiểm soát với các mô hình cơ sở, báo cáo
  chất lượng, độ phù hợp cá nhân hoá, độ an toàn tái sử dụng, độ trễ và chi phí bằng
  phương pháp kiểm định thống kê phù hợp.
\end{enumerate}
